\documentclass[11pt]{article}

\usepackage[a4paper,margin=1in]{geometry}
\usepackage{amsmath,amssymb}
\usepackage[hidelinks]{hyperref}
\usepackage{enumitem}

\title{Review Findings for \texttt{rational\_cut\_and\_project\_gap\_periods.tex}}
\author{Codex Review Notes}
\date{May 7, 2026}

\begin{document}

\maketitle

\section*{Executive Summary}

I reviewed
\texttt{LaTeX/rational\_cut\_and\_project\_gap\_periods.tex}
with an emphasis on mathematical correctness, arXiv readiness,
reproducibility, and consistency with the Lean and CSV artifacts in the
repository.

The main rational proof line appears mathematically sound: the strip
condition reduces correctly to an interval of admissible invariants
\[
  r = \alpha x - \beta y,
\]
each invariant class gives one arithmetic progression in the
\(\tilde p\)-coordinate, and the multiset and set period formulas follow
from the residue-distribution dichotomy modulo
\[
  D = \alpha^2+\beta^2.
\]
I did not find a mathematical flaw in the main rational period theorem.

The most important issues before arXiv submission are not in the core
formula, but in reproducibility and presentation.  In particular, the
empirical verification claims in the paper are currently too strong for the
CSV data in the repository.

\section*{Publication-Critical Findings}

\begin{enumerate}[label=\textbf{F\arabic*.},leftmargin=*]

\item \textbf{The empirical verification claim is currently false as written.}

In
\[
\texttt{LaTeX/rational\_cut\_and\_project\_gap\_periods.tex:414--438}
\]
and again at
\[
\texttt{LaTeX/rational\_cut\_and\_project\_gap\_periods.tex:1088--1090},
\]
the manuscript states that all \(51{,}729\) CSV cases match both the multiset
formula and the set formula with zero failures.

For the files currently present in \texttt{tests/}, the multiset formula does
match all \(51{,}729\) data rows.  The set formula, however, matches only
\(51{,}686\) rows.  There are \(43\) rows in
\[
\texttt{tests/multiset\_and\_set\_new\_find\_patterns\_51012\_lines.csv}
\]
where \texttt{period\_set} is a negative error code rather than a period.
For example, CSV line 3345 has
\[
  N=35{,}103{,}140,\qquad D=7{,}720{,}913{,}978,
\]
so the theorem predicts
\[
  \lambda_{\mathrm{set}} = N = 35{,}103{,}140,
\]
but the stored \texttt{period\_set} value is \texttt{-2}.

\textbf{Recommendation.}  Either regenerate the affected set-period column
so that all rows contain genuine periods, or rewrite the empirical remark to
separate proved theorem checks from incomplete/failed numerical set-period
computations.  For arXiv, the cleanest option is to cite only the verified
rows, or to replace the CSV with a fully theorem-consistent file.

\item \textbf{The verification script is out of sync with the current CSV
format.}

The script
\[
\texttt{src/python/verify\_paper.py}
\]
expects five columns
\[
\texttt{o\_n,o\_d,a\_n,a\_d,period},
\]
but the current main CSV files contain six columns:
\[
\texttt{o\_n,o\_d,a\_n,a\_d,period\_multiset,period\_set}.
\]
Running the script on
\[
\texttt{tests/multiset\_and\_set\_new\_find\_patterns\_51012\_lines.csv}
\]
therefore fails with
\[
\texttt{ValueError: too many values to unpack (expected 5)}.
\]

\textbf{Recommendation.}  Update the verification script so that it checks
both formulas against the six-column CSV format, reports negative
\texttt{period\_set} values explicitly, and prints the exact aggregate counts
used in the paper.

\end{enumerate}

\section*{Mathematical Precision Findings}

\begin{enumerate}[label=\textbf{M\arabic*.},leftmargin=*]

\item \textbf{The residue-bijection corollary should include a proof.}

At
\[
\texttt{LaTeX/rational\_cut\_and\_project\_gap\_periods.tex:605--609},
\]
the residue-bijection corollary states that the map
\[
  r \mapsto c_r = -\alpha\beta^{-1}r \pmod D
\]
is a bijection on \(\mathbb Z/D\mathbb Z\), but no proof is given.

The result is correct.  A short proof would improve the paper:
\(\beta^{-1}\) exists because \(\gcd(\beta,D)=1\), and
\(-\alpha\beta^{-1}\) is a unit modulo \(D\) because
\(\gcd(\alpha,D)=1\).  Multiplication by a unit is a bijection.

\item \textbf{Clarify that \(D\) is a common difference in the
\(\tilde p\)-coordinate.}

In the introduction, around
\[
\texttt{LaTeX/rational\_cut\_and\_project\_gap\_periods.tex:135--137},
\]
the manuscript says that \(D\) is the common difference of the arithmetic
progressions obtained after projection.

Strictly speaking, \(D\) is the common difference in the integer coordinate
\[
  \tilde p = \beta x + \alpha y.
\]
The Euclidean distance along the projected line is obtained from
\(\tilde p\) by multiplication by a positive scale factor.  The later text
mostly makes this clear, but the introductory statement should be sharpened.

\item \textbf{The irrational-case setup should explicitly justify
unboundedness.}

Around
\[
\texttt{LaTeX/rational\_cut\_and\_project\_gap\_periods.tex:1217--1226},
\]
the manuscript says that the irrational projected set is locally finite and
unbounded in both directions.  The local finiteness argument is clear.  The
unboundedness assertion should be justified explicitly, for example from the
density of \(s(\mathbb Z^2)\) and the fact that the window
\[
  W=[-a\omega,\omega]
\]
has nonempty interior when \(\omega>0\).

The later proof of aperiodicity appears correct, but this preliminary
well-definedness step would benefit from one or two added sentences.

\item \textbf{The main minimality argument is sound, but one bridge could be
made more explicit.}

In the proof of the generic-period proposition, especially around
\[
\texttt{LaTeX/rational\_cut\_and\_project\_gap\_periods.tex:787--839},
\]
the key step is that a smaller gap period gives a translation symmetry of
the entire projected multiset, hence of the residue multiplicity function.
This is correct.  For maximum readability, it would help to add one sentence
spelling out that translation by an integer \(\sigma\) sends residue class
\(c\) modulo \(D\) to \(c+\sigma\), so multiset invariance forces
\[
  \mu(c+\sigma)=\mu(c)
\]
for every residue class.

\end{enumerate}

\section*{Lean and Repository Findings}

\begin{enumerate}[label=\textbf{L\arabic*.},leftmargin=*]

\item \textbf{The current Lean file builds successfully.}

Running
\[
\texttt{lake build}
\]
inside \texttt{Lean/CutAndProject} completed successfully.  The current file
\[
\texttt{Lean/CutAndProject/CutAndProject/Basic.lean}
\]
contains no occurrences of \texttt{sorry}, \texttt{admit}, or \texttt{axiom}.

\item \textbf{Stored build logs are stale.}

The repository files
\[
\texttt{Lean/CutAndProject/build\_output.txt}
\quad\text{and}\quad
\texttt{Lean/CutAndProject/build\_out.txt}
\]
contain old build failures and old \texttt{sorry} warnings.  These logs are
inconsistent with the current successful build.

\textbf{Recommendation.}  Delete these logs before release, or replace them
with current successful build output.

\item \textbf{Small path typo in the Lean correspondence table.}

At
\[
\texttt{LaTeX/rational\_cut\_and\_project\_gap\_periods.tex:1444--1447},
\]
the table caption refers to
\[
\texttt{LEAN/CutAndProject/CutAndProject/Basic.lean}.
\]
The repository path is
\[
\texttt{Lean/CutAndProject/CutAndProject/Basic.lean}.
\]

\end{enumerate}

\section*{LaTeX and arXiv Readiness}

\begin{enumerate}[label=\textbf{A\arabic*.},leftmargin=*]

\item \textbf{The paper compiles successfully.}

Running
\[
\texttt{latexmk -g -pdf -interaction=nonstopmode -halt-on-error}
\]
on the manuscript completed successfully.  There were no undefined
references or undefined citations.

\item \textbf{There are several harmless but fixable hyperref warnings.}

The log reports warnings of the form
\[
\texttt{Token not allowed in a PDF string}
\]
for headings containing mathematics, for example
\(\omega=0\), \(D\nmid N\), and \(D\mid N\).

\textbf{Recommendation.}  Use \verb|\texorpdfstring| in
section or subsection headings containing mathematical notation.

\item \textbf{Underfull boxes are present but not publication-critical.}

The compilation log contains several underfull \verb|\hbox|
warnings, mostly caused by long \verb|\texttt| identifiers
and narrow table/caption text.  I did not see overfull boxes.  These warnings
are not mathematically significant, but the Lean correspondence table and
long file names could be visually polished before submission.

\end{enumerate}

\section*{Recommended Action List}

\begin{enumerate}[leftmargin=*]

\item Fix the empirical verification claim before arXiv submission.  This is
the only issue I regard as publication-critical.

\item Update \texttt{src/python/verify\_paper.py} to the current six-column
CSV format and use it to regenerate the exact numbers quoted in the paper.

\item Add a short proof of the residue-bijection corollary.

\item Clarify that \(D\) is the common difference in the \(\tilde p\)-coordinate,
not directly the Euclidean projected distance.

\item Add an explicit sentence proving unboundedness of the irrational
projected set.

\item Remove or refresh stale Lean build logs.

\item Fix the path capitalization \texttt{LEAN} \(\to\) \texttt{Lean}.

\item Optionally clean the hyperref warnings caused by mathematics in PDF
bookmarks.

\end{enumerate}

\section*{Overall Assessment}

The rational main theorem and its set-valued companion appear strong and
mathematically coherent.  The proof strategy is well chosen: it isolates the
geometry in the strip-to-residue reduction and then proves the period
formulas through finite residue-counting and cyclic-interval symmetry.

After correcting the empirical-verification passage and the small
reproducibility issues above, the manuscript should be much closer to a
high-quality arXiv submission.

\end{document}
